%1. La memoria incluirá una portada portada normalizada con la siguiente información: título en castellano, título en inglés, autores, profesor director, codirector si es el caso, curso académico e identificación de la asignatura (Trabajo de fin de grado del Grado en - nombre del grado correspondiente-, Facultad de Informática, Universidad Complutense de Madrid). Los datos referentes al título y director (y codirector en su caso) deben corresponder a los publicados en web de TFG, según lo indicado en el punto 6 de la sección III de esta normativa. 
% 2. en cada apartado
% 3. La memoria constará de un mínimo de 25 páginas para los proyectos realizados por unúnico estudiant

\documentclass[12pt , a4paper]{article}
\usepackage[spanish, english]{babel} 	% language
\usepackage{tocloft}					% table of contents

% for sections
\renewcommand{\cftsecleader}{\cftdotfill{\cftdotsep}} 

% images
\usepackage{graphicx}
\graphicspath{ {./images/} }

% text columns
\usepackage{multicol}
\setlength{\columnsep}{1cm}


\title{\textbf{Sampler plugin for digital audio workstations}}

%\title{\textbf{Sampler plugin for digital audio workstations\\---\\ Instrumento virtual basado en samples para estaciones de trabajo de audio digital}}

%\author{Juan Chozas Sumbera}
\date{\vspace{-10ex}}

\begin{document}
	\maketitle
	\begin{center}
		\textbf{Trabajo de Fin de Grado en Ingeniería Informática}\\
		
		~\newline
		\textbf{\large{Juan Chozas Sumbera}}
		
		~\newline
		Dirigido por
		
		\textbf{Jaime Sánchez Hernández\\
		Miguel Gómez-Zamalloa Gil}
		
		
		~\newline
		\includegraphics[width=0.5\textwidth]{logo_UCM.png}\\
		~\newline
		
		\textbf{\large{UNIVERSIDAD COMPLUTENSE DE MADRID}}\\
		
		\textbf{FACULTAD DE INFORMÁTICA}
	\end{center}

	
	
	% 2. La memoria debe incluir la descripción detallada de la propuesta hardware/software realizada y ha de contener: 
	% a. un índice, 
	% \tableofcontents % despues del resumen y antes de la intro
	\newpage

	
	% b. un resumen y una lista de no más de 10 palabras clave para su búsqueda bibliográfica, ambos en castellano e inglés, 
	\newpage
	\huge
	\textbf{Abstract}\\

	\normalsize
	When it comes to modern music composition, Musical Instrument Digital Interface (MIDI from now on) protocols provide grounds for an alternative to \textit{analog} instruments. This standard erects a bridge for communication between digital instruments and computers, providing the world of digital music production with the tangible interfaces of normal instruments. The MIDI standard provides, among many other features, a communications protocol that encodes music events. This manifests in the form of messages which describe data such as the pressing, releasing, and velocity of musical notes. A digital audio workstation (DAW from here on out) is a software environment for music production. These workstations revolve around MIDI, and can be extended with plugins in the form of effects or instruments. The former normally have a sound-based input-output relation with the DAW, meanwhile instruments will take MIDI messages and produce a sound as an output. For example, a simple plugin could be a sine wave based instrument. The frequency of the wave will depend on the note being held down, and the force applied to the key will dictate the amplitude of the wave. This plugin, although simple, could be considered a synthesizer.
		
	\par
	A sampler is a synthesizer that uses samples rather than oscillators to produce sound.  The sources of the sounds it produces are samples, which can be described as clips of audio that are stored in a digital format within the sampler's memory. You're likely to find that hardware and software samplers offer sets of "stock" samples that are loaded by default. The option to load samples from an external source such as an SD card is also common, just like recording external input from a recording device such as a microphone. As is the case with non sample based synthesizers, there are several ways that the samples can be manipulated. Once samples are loaded, the sampler can make the notes on a keyboard trigger the samples so users can play them to compose music. This sampler, in particular, takes the form of a plugin for digital audio workstations. The DAW routes MIDI messages describing the notes being played into the plugin, which in turn processes the input to produce a stream of audio that is returned to the host as sound.\\
	

	\large
	\textbf{Key words}\\
	
	\vspace{-1em}
	\normalsize	
	\noindent sound, sampler, sampling, instrument, MIDI, DAW, plugin, VST, synthesizer\\

	\newpage
	\huge
	\textbf{Resumen}\\
	
	\normalsize
	Hoy en día el mundo de la producción musical dispone del Musical Instrument Digital Interface (MIDI de ahora en adelante), una serie de protocolos que permite el uso de alternativas a instrumentos \textit{analógicos}. El MIDI hace de puente entre ordenadores e instrumentos digitales, dando lugar a una interfaz tangible en el mundo de la producción digital. Éste estándar ofrece, entre otras cosas, un protocolo de comunicación que codifica información musical. Éste consiste en mensajes que describen eventos cómo lo pueden ser la pulsación de una nota, el soltar una nota, o la fuerza con la que se pulsa una nota. Una estación de audio digital (DAW, abreviado en inglés de Digital Audio Workstation) es una clase de software que presenta un entorno digital para la producción musical en el que el estándar MIDI es fundamental. Los DAWs pueden ser extendidos en cuanto a efectos e instrumentos mediante plugins. Los efectos suelen recibir un canal de sonido que, una vez procesado, se devuelve al DAW. En cambio, el sonido devuelto por los instrumentos depende de los mensajes MIDI que recibe como entrada. Por ejemplo, un instrumento de ondas sinusoidales sería un plugin sencillo. La frecuencia de las ondas dependería de la nota accionada y su amplutud se vería afectada por la fuerza con la que se presionan las teclas. Este sencillo plugin se consideraría un sintetizador, a pesar de su sencillez.

	\par
	Un sampler es un sintetizador que produce sonido mediante samples en lugar de generar sonidos con osciladores. Las fuentes del sonido que produce son samples: muestras de sonido grabado en forma digital que se almacena en en la memoria del sampler. Muchos samplers vienen con una serie de samples cargados por defecto, pero siempre ofrecen la opción de cargar samples distintos. Esto se puede hacer a través de una unidad de almacenamiento, como una tarjeta SD, o con una señal externa usando un dispositivo con capacidad de grabación como la que tiene un micrófono. Igual que con los sintetizadores, hay múltiples formas de manipular los samples para producir sonidos personalizados. Una vez cargados los samples, se pueden asociar a las notas de un teclado para accionar su reproducción y permitir que el usuario componga música. Este sampler está implementado como plugin para estaciones de audio digital, y funciona tomando como entrada disparadores de la estación de audio digital para devolver las notas recibidas en forma de sonido.\\
	
	\large
	\textbf{Palabras clave}\\
	
	\vspace{-1em}
	\normalsize	
	\noindent sonido, muestra, muestreo, instrumento, MIDI, DAW, plugin, VST, sintetizador
	
	
	
	\newpage
	% a. un índice, 
	\tableofcontents


	% c. una introducción con los antecedentes, objetivos y plan de trabajo, 
	\newpage
	\section{Introduction}
	\subsection{History of sampling}
	
	The technique of sampling dates back to the 1960s when recordings were captured on tape. Thanks to a hardware design that made contact between the playback head and different sections of a moving tape, musicians could play the distinct recordings on keyboards of instruments such as the Mellotron. During the 1980s, the popularity of drum machines increased significantly. Many drum machines were sample-based, i.e. they created sounds using digitally stored samples. The alternative was to synthesize sound in an analog fashion. In 1988, the first Music Production Center (MPC) by Akai was made available to the public: a sample-based music workstation that was capable of arranging samples of all lengths in its sequencer to produce full fledged music tracks without the need for additional instruments or hardware. In a recipe for music, an MPC could be the only ingredient, reaching a vast number of musicians because of its affordability in comparison to previous means of music production. It is a tool that gave artists a new way to create music, a technique that has been a foundation to several genres and highly influential to music as a whole.
	\par	
	In a world of highly complex information systems, digital music production has been developed thoroughly. One can have an extensive range of digital synthesizers, samplers, and effect chains all working simultaneously in computer software. A DAW provides a playground that routes audio between racks of components and effects. With support for plugins, such as instruments or effects, the software's capabilities can be extended to broaden a musician's repertoire for music production. 

	\newpage
	\subsection{Objectives}
	Most DAWs have built in samplers or similar tools that allow a user to sample sound and directly manipulate it so as to employ sampling techniques. In most cases, these tools are more than capable of providing the means to translate these techniques into action, however, depending on the DAW, you may be hindered due to the program's presentation. Whether it be the user interface, the imposed workflow, or the minimum amount of steps required to reach your sampling goals, it is likely that the process involves inconveniences. Personally, I find that there are a minimum of three steps to achieve a usable sampler configuration. First, a main sample must be loaded. Second, the process of \textit{chopping} the sample, where a set of subsections are created within the sample, creating sub-samples otherwise known as \textit{chops}. Third comes the assignment of chops to MIDI notes, in other words, the mapping of sample subsection playback to a controller's keys. These steps are a prerequisite to playing chops on the controller. At this stage, the user can experience the instrument and judge whether it is necessary to take a step back and make adjustments, for example, shifting the start time of a chop, creating/deleting a chop, or moving the trigger note to another of higher convenience. All of this implies that the design of a sampler instrument is an iterative process, and that the user may keep reforming it in several ways to fit the necessities of the creative process in musical composition.
	\par	
	The aim here is to reduce the amount of user interaction required to translate an idea into sound. For the simple implications of sampling techniques, the means by which you achieve a minimum setup with which you can play around and manifest ideas should be straightforward. With this project the intention is to strip the process of whatever element that is not essential to it. By downsizing in the features that a DAW would offer, I aspire to build a plugin strictly optimized for sampling. This means that the user will be able to load, chop, and play samples with speed and ease in a reduced sequence of actions. In addition, the user should also be able to make fast corrections and tweaks to further develop the instrument. 
	
	% d. resultados y discusión crítica y razonada de los mismos, con sus conclusiones,
	\newpage
	\section{Results}
	
	\subsection{Conclusion}
	 
	\newpage
	% e. bibliografía.
	\section{Bibliography}
	
\end{document}